% Secciones 4--7 del delta gravitatorio íntegro.
\section{Incidence hexade--octade, paramètre de Barbero--Immirzi et spectre d'aire}
\label{sec:l9s102:barbero-area}

La dérivation surfacique reçoit comme antécédents l'inclusion d'incidence
\(6\to8\), les cardinaux TPK \(90/120\), le défaut normalisé
\(-1/12\) et les canaux spectraux \(q_\pm\). La fonctionnelle
hexade--octade détermine le paramètre de Barbero--Immirzi et la
normalisation de l'opérateur d'aire. Cet emplacement contient la dérivation
HMT et le spectre ; la réalisation Palatini--Cartan--Holst occupe un emplacement
ultérieur dans la mécanique dimensionnelle.

Soient \(H\subset O\), avec \(|H|=6\) et \(|O|=8\), et conservons le couple
marqué de l'hexade dans l'inclusion
\begin{equation}
 \iota_{6,8}:\mathcal F_6=H\times\binom H2
 \lhook\joinrel\longrightarrow
 \mathcal F_8=O\times\binom H2.
 \label{eq:l9s102:barbero-inclusion-incidencia}
\end{equation}
Le comptage interne donne
\begin{equation}
 |\mathcal F_6|=6\binom62=90,
 \qquad |\mathcal F_8|=8\binom62=120,
 \qquad
 \frac{90-120}{24\binom62}=-\frac1{12}.
 \label{eq:l9s102:barbero-conteos}
\end{equation}
Une fois \(\pi_{\rm HMT}\), l'angle \(A\) et le
contre-angle unique \(C^*\) publiés, les deux canaux sont
\begin{equation}
 q_\pm=\exp\!\left[-\frac{\pi_{\rm HMT}}{180}(A\pm C^*)\right],
 \qquad 0<q_+<q_-<1,
 \label{eq:l9s102:barbero-qpm}
\end{equation}
et leurs séries de retour orienté sont
\begin{equation}
 S_{90}(q)=\frac{q^{90}}{1-q^{270}},\qquad
 S_{120}(q)=\frac{q^{120}}{1-q^{360}},\qquad
 \Delta S_m=S_m(q_-)-S_m(q_+).
 \label{eq:l9s102:barbero-series}
\end{equation}

L'ordre causal des coefficients est fixé dans le registre temporel du TPK,
avant d'évaluer le coefficient du pôle. Soit \(C_{12}=\mathbb Z/12\mathbb Z\) la roue
dodécaphasique et soit \(\partial_{\rm ret}\) son unique couture orientée de retour.
Dans le module libre des événements d'incidence, le tour fondamental est
\begin{equation}
 c_{12}^{+}
 =\sum_{j\in C_{12}}[j,\mathcal F_6]
  -[\partial_{\rm ret},\mathcal F_8].
 \label{eq:l9s102:barbero-cadena-dodecafasica}
\end{equation}
Le signe du second terme est le signe de bord. L'involution
bidirectionnelle inverse l'orientation complète,
\(c_{12}^{-}=-c_{12}^{+}\), sans modifier les multiplicités absolues des
deux types d'incidence. Définissons le lecteur de retour sur les
générateurs par
\begin{equation}
 \mathscr R([j,\mathcal F_6])=S_{90},
 \qquad
 \mathscr R([\partial_{\rm ret},\mathcal F_8])=S_{120}.
 \label{eq:l9s102:barbero-lector-retorno}
\end{equation}
La covariance de phase rend la première image constante sur les douze
secteurs ; la seconde agit uniquement sur la couture élargie. Par linéarité,
\begin{equation}
 \boxed{\mathscr R(c_{12}^{+})=12S_{90}-S_{120}.}
 \label{eq:l9s102:barbero-peso-generado}
\end{equation}
Ainsi, le poids \(12:-1\) est l'image du tour fondamental : douze secteurs et
un bord orienté. Il ne provient ni d'une imposition préalable de \(1/24\), ni d'une
comparaison décimale, ni de la valeur conventionnelle de Barbero--Immirzi.

\Needspace{12\baselineskip}
\begin{theorem}[Détermination dodécaphasique de la fonctionnelle de Barbero--Immirzi]
\label{thm:l9s102:barbero-funcional}
La chaîne fondamentale \eqref{eq:l9s102:barbero-cadena-dodecafasica} détermine
le couple orienté \((a,b)=(12,1)\). Le coefficient de \((1-q)^{-1}\) de son image est
alors \(1/24\), comme sortie exacte de la fonctionnelle déjà construite. Par conséquent,
après évaluation des canaux conjugués \(q_\pm\), la sortie adimensionnelle est
\begin{equation}
 \boxed{
 \gamma_{\rm HMT}=\Gamma_{6\to8}
 =\frac{\pi_{\rm HMT}AC^*}{180}
   +12\Delta S_{90}-\Delta S_{120}.}
 \label{eq:l9s102:barbero-gamma}
\end{equation}
Cette formule reçoit le couple angulaire déjà généré et n'utilise pas en entrée
la valeur conventionnelle du paramètre de Barbero--Immirzi.
\end{theorem}

\begin{proof}
La somme de \eqref{eq:l9s102:barbero-cadena-dodecafasica} contient exactement
les douze secteurs de \(C_{12}\), tandis que la couture de retour apparaît une
seule fois et avec une orientation négative. Le lecteur
\eqref{eq:l9s102:barbero-lector-retorno} produit donc
\eqref{eq:l9s102:barbero-peso-generado}, avant tout calcul de pôle. Pour
chaque \(m>0\),
\[
 p_1(S_m)
 =\lim_{q\to1}(1-q)\frac{q^m}{1-q^{3m}}
 =\frac1{3m}.
\]
En conséquence,
\[
 p_1\bigl(\mathscr R(c_{12}^{+})\bigr)
 =\frac{12}{270}-\frac1{360}
 =\frac{48-3}{1080}=\frac1{24}.
\]
Le coefficient \(p_1\) est rapporté à \((1-q)^{-1}\).
Dans la coordonnée complexe \(q-1\),
\[
 \operatorname*{Res}_{q=1}S_m(q)=-\frac1{3m},
 \qquad
 \operatorname*{Res}_{q=1}\mathscr R(c_{12}^{+})(q)=-\frac1{24}.
\]
L'équation diophantienne ultérieure
\[
 \frac a{270}-\frac b{360}=\frac1{24}
 \quad\Longleftrightarrow\quad 4a-3b=45
\]
possède les solutions entières positives
\((a,b)=(12+3t,1+4t)\), \(t\ge0\) ; l'unique solution minimale est
\((12,1)\). Ce calcul certifie arithmétiquement le couple que l'orbite
fondamentale avait déjà engendré : tout \(t>0\) altère la multiplicité des
secteurs ou de la couture et cesse de représenter un tour du TPK.
Enfin, substituer \(q_\pm\) dans l'image orientée et ajouter la composante
angulaire préalablement générée prouve \eqref{eq:l9s102:barbero-gamma}.
\end{proof}

La séquence causale est fixée sans échange des rôles :
\begin{equation}
\begin{aligned}
 (\mathcal F_6\hookrightarrow\mathcal F_8,C_{12},\partial_{\rm ret})
 &\longrightarrow (90,120,-1/12,c_{12}^{+})\\
 &\longrightarrow \mathscr R(c_{12}^{+})=12S_{90}-S_{120},\\
 \mathscr R(c_{12}^{+})
 &\longrightarrow p_1\bigl(\mathscr R(c_{12}^{+})\bigr)=\frac1{24},\\
 \bigl(\mathscr R(c_{12}^{+}),A,C^*,q_\pm\bigr)
 &\longrightarrow \Gamma_{6\to8}=\gamma_{\rm HMT}
 \longrightarrow \widehat{\mathcal A}_{\rm phys}.
\end{aligned}
 \label{eq:l9s102:barbero-cadena-completa}
\end{equation}
Le couple angulaire et ses canaux proviennent de la construction précédente.
L'évaluation du coefficient du pôle constitue une vérification
de la fonctionnelle de retour ; l'évaluation conjuguée emploie les
coordonnées angulaires préalablement générées.
La réalisation ultérieure dans l'action de Holst reconnaît la même section
adimensionnelle ; elle n'intervient pas dans sa sélection.

Pour les deux secteurs de bord, soient
\begin{equation}
 N_m=\sum_{e\in\partial A_m}N_e,
 \qquad m\in\{90,120\},
 \qquad \operatorname{Spec}(N_e)=\{0,1,2\}.
 \label{eq:l9s102:n90-n120}
\end{equation}
L'opérateur physique d'aire s'écrit
\begin{equation}
 \boxed{
 \frac{\widehat{\mathcal A}_{\mathrm{phys}}}{\ell_P^2}
 =\gamma_{\mathrm{HMT}}a_0(N_{90}+N_{120}).}
 \label{eq:l9s102:area-operador}
\end{equation}
Sur \(36\) liens qutrit,
\begin{equation}
 \operatorname{Spec}(N_{90}+N_{120})=\{0,1,\ldots,72\},
 \label{eq:l9s102:area-numero-spectrum}
\end{equation}
et la multiplicité de la valeur propre \(n\) est le coefficient de \(z^n\) dans
\((1+z+z^2)^{36}\). Il en résulte
\begin{equation}
 \operatorname{Spec}\!\left(
 \frac{\widehat{\mathcal A}_{\mathrm{phys}}}{\ell_P^2}\right)
 =\{n\gamma_{\mathrm{HMT}}a_0:0\le n\le72\}.
 \label{eq:l9s102:area-spectrum}
\end{equation}

\section{Géométrie totale et mémoire de provenance}
\label{sec:l9s102:area-hamiltoniano}

Le hamiltonien modulaire du bord est
\begin{equation}
 \boxed{
 K_A=-\log\rho_A
 =cI+\Delta_{\mathrm{mod}}(90N_{90}+120N_{120}).}
 \label{eq:l9s102:hamiltoniano-modular}
\end{equation}
Soient
\begin{equation}
 N=N_{90}+N_{120},
 \qquad
 D=90N_{90}+120N_{120}.
 \label{eq:l9s102:n-d}
\end{equation}
La matrice de passage
\(\left(\begin{smallmatrix}1&1\\90&120\end{smallmatrix}\right)\)
a pour déterminant \(30\), et son inverse reconstruit les secteurs :
\begin{equation}
 \boxed{
 N_{90}=4N-\frac{D}{30},
 \qquad
 N_{120}=\frac{D}{30}-3N.}
 \label{eq:l9s102:reconstruccion-sectores}
\end{equation}
L'aire publie le nombre total d'excitations ; le hamiltonien modulaire
publie leur secteur de provenance. Le couple conjoint reconstruit l'occupation
complète du bord.

\section{État bicouche et information orientée}
\label{sec:l9s102:estado-bicapa}

Soit \(R=R^*=R^{-1}\) l'involution des feuillets et
\(y=C^*_{\mathrm{rad}}\). L'état normalisé est
\begin{equation}
 \boxed{
 \rho_y=\frac{e^{-yR}}{2\cosh y}.}
 \label{eq:l9s102:rho-y}
\end{equation}
Son hamiltonien modulaire et son entropie sont
\begin{equation}
 K_y=-\log\rho_y=\log(2\cosh y)I+yR,
 \label{eq:l9s102:ky}
\end{equation}
\begin{equation}
 S(\rho_y)=\log(2\cosh y)-y\tanh y.
 \label{eq:l9s102:entropia-y}
\end{equation}
L'inversion \(y\mapsto-y\) laisse l'entropie fixe et produit la distance
informationnelle orientée
\begin{equation}
 \boxed{
 D(\rho_y\Vert\rho_{-y})=2y\tanh y.}
 \label{eq:l9s102:entropia-relativa-hojas}
\end{equation}
La composante paire mesure la distribution spectrale ; la composante impaire mesure
le coût informationnel de l'inversion de l'orientation.

\section{Quantum informationnel d'énergie et cellule gravitationnelle de demi-action}
\label{sec:l9s102:horizonte-media-accion}

La périodicité de \(108\) états détermine
\begin{equation}
 t_P=\frac{54}{\pi}t_0,
 \qquad
 E_P=\frac{\hbar}{t_P}=\frac{h}{108t_0}.
 \label{eq:l9s102:ep-tp}
\end{equation}
La réalisation thermique de l'horloge tritique satisfait
\begin{equation}
 \boxed{
 E_P=k_B\Theta_{\mathrm{clk}}\log3.}
 \label{eq:l9s102:planck-trit}
\end{equation}
Pour l'horizon de rayon \(R\),
\begin{equation}
 E_\Lambda(R)=\frac{c^4R}{2G},
 \qquad
 E_\Lambda(R)=T_H(R)S_H(R).
 \label{eq:l9s102:energia-horizonte}
\end{equation}
Dans la cellule de Planck,
\begin{equation}
 \boxed{
 E_\Lambda
 =T_HS_H
 =\frac{E_P}{2}
 =\frac{k_B\Theta_{\mathrm{clk}}\log3}{2}.}
 \label{eq:l9s102:confluencia-horizonte}
\end{equation}
Par conséquent,
\begin{equation}
 \boxed{
 E_\Lambda t_P=\frac{\hbar}{2},
 \qquad
 E_\Lambda T_{108}=\frac h2.}
 \label{eq:l9s102:media-accion}
\end{equation}
L'énergie du vide, l'entropie surfacique, la température chronologique et
l'action partagent une même cellule de confluence avec des domaines typés.

Lorsque \(S_H/k_B=\pi\), la multiplicité effective satisfait
\begin{equation}
 \Omega_{\mathrm{eff}}=e^{S_H/k_B}=e^\pi.
 \label{eq:l9s102:omega-e-pi}
\end{equation}
L'opérateur hyperbolique d'Euler possède
\begin{equation}
 \operatorname{Spec}(e^{\pi H})=\{e^\pi,e^{-\pi}\}.
 \label{eq:l9s102:e-pi-spectrum}
\end{equation}
L'égalité des scalaires établit la confluence spectrale ;
l'opérateur d'entrelacement entre multiplicité microscopique d'horizon et dynamique de
Perron relève de son emplacement physique spécifique.
